\section{A linear-degree local representation of deterministic machines}
\label{btc:mach:section}

The connection with automata requires more than a matching representation
of a completed configuration graph. Each input symbol must determine its
own diagram, independently of its neighbors and its position. We adapt the
candidate-slot construction of \cite{OpenAILiveness2026}, retaining only
crossing transitions that occur somewhere in the finite transition table.
This changes its quadratic degree bound to a linear bound for every fixed
alphabet. The underlying traversal of a configuration tree is related to
the halting and reversible-simulation methods of
\cite{SipserHalting1980,LangeMcKenzieTapp2000}; all traversal and locality
claims needed here are proved below.

\subsection{Machine convention and statement}

Let $\Sigma$ be an alphabet of size $a$. The scanned input is
$\Lmark w\Rmark$, where the two endmarkers are distinct and are not in
$\Sigma$. An $s$-state deterministic two-way finite automaton has state
set $Q$, initial state $q_0$, accepting set $A\subseteq Q$, and a partial
function
\[
 \delta:Q\times(\Sigma\cup\{\Lmark,\Rmark\})
 \rightharpoonup Q\times\{-1,0,1\}.
\]
Moves at endmarkers never leave the marked tape. The initial head position
is the left marker. Acceptance means that an accepting state occurs at
some finite time, including time zero. An undefined transition halts; an
infinite run without an accepting state rejects. Initial and accepting
states are included in $s$. Stay moves, partial transition tables, and
rejecting cycles are allowed throughout.

\begin{theorem}[Sparse-slot local representation]\label{btc:mach:representation}
For every machine in this convention, there are an even integer $k$ and
a map
\[
 \eta:\Sigma\cup\{\Lmark,\Rmark\}\longrightarrow\Br{k}
\]
such that
\begin{equation}\label{btc:mach:degree-bound}
 2\le k\le 2(a+2)s+2.
\end{equation}
Designate ports $1,2$ on each side as test ports. For every
$w=a_1\cdots a_\ell$, the machine accepts $w$ exactly when
\begin{equation}\label{btc:mach:whole-diagram}
 \eta(\Lmark)\eta(a_1)\cdots\eta(a_\ell)\eta(\Rmark)
\end{equation}
pairs a left test port with a right test port. In that case both left
test ports are paired with right test ports. Each diagram depends only
on its displayed symbol and the machine. In particular, for
$\Sigma=\bits$ one has $k\le8s+2$.
\end{theorem}

\subsection{Acceptance becomes reachability of one sink}

Add a new state $f$ and define an auxiliary partial successor rule by
\begin{equation}\label{btc:mach:normalization}
 \widehat\delta(q,b)=
 \begin{cases}
 (f,0),&q\in A,\\
 \delta(q,b),&q\in Q\setminus A,\\
 (f,1),&q=f,\ b\ne\Rmark,\\
 \text{undefined},&q=f,\ b=\Rmark.
 \end{cases}
\end{equation}
The accepting-state rule takes precedence even when the original entry
was undefined. For a fixed word of length $\ell$, let
$c_0=(0,q_0)$ and $t=(\ell+1,f)$. The original machine accepts exactly
when this new successor rule reaches $t$ from $c_0$. Until an accepting
state first occurs, the runs agree. Such a state enters $f$, which sweeps
to the right marker and stops. Conversely, the first entry into $f$
must come from an old accepting state. This includes acceptance initially
and acceptance on either marker. There is no additional copy of the
initial state.

\begin{lemma}[The component of a functional sink]\label{btc:mach:tree}
Let $\sigma:V\rightharpoonup V$ be a partial function on a finite set,
and suppose $\sigma(t)$ is undefined. Make an undirected multigraph with
one edge for each defined successor, retaining loops and parallel edges.
The component of $t$ consists precisely of the vertices whose forward
successor paths reach $t$, and this component is a tree.
\end{lemma}

\begin{proof}
Write $B$ for the set of vertices reaching $t$, allowing zero steps.
For every successor $u\mapsto v$, one has $u\in B$ if and only if
$v\in B$: prepend the edge in one direction, and remove the first
step of the unique successor path in the other. Here $u\ne t$, because
$t$ has no successor. Membership in $B$ is therefore constant along
undirected edges. A forward path is also an undirected path, so $B$ is
exactly the component of $t$.

Every vertex of that component other than $t$ has exactly one outgoing
edge, and all these edges stay in the component. If the component has
$v$ vertices, it consequently has $v-1$ edges, counted with multiplicity.
It is connected and hence is a tree, with neither loops nor parallel
edges.
\end{proof}

Apply the lemma to all configurations of \eqref{btc:mach:normalization}.
Call the resulting undirected successor multigraph $G$. Its component
containing $t$ is a tree, and contains $c_0$ precisely on accepted
inputs. Other components may contain nonaccepting cycles. They need no
modification.

\subsection{A sparse set of boundary slots}

Put $Q'=Q\cup\{f\}$. The fixed set of crossing candidates is
\begin{equation}\label{btc:mach:slot-set}
 C=\{(d,p,q):d\in\{-1,1\},\ p,q\in Q',\
       \widehat\delta(p,b)=(q,d)\text{ for some }b\},
 \qquad J=|C|.
\end{equation}
Each old state supplies at most $a+2$ candidates, one per scanned symbol.
The new state supplies just the single candidate $(1,f,f)$, since all
its defined crossing rules are identical. Therefore
\begin{equation}\label{btc:mach:slot-count}
 1\le J\le(a+2)s+1,\qquad k=2J\le2(a+2)s+2.
\end{equation}
In fact the first upper bound can be improved to
$(a+2)(s-|A|)+1$, because old accepting states now have only stay rules.
The important point is that the candidate set is a union over local
transition-table entries, not over all pairs of states.

Fix an enumeration of $C$. At every boundary between neighboring cells,
reserve one edge slot for each $(d,p,q)\in C$. The direction $d$
determines which cell is the source: the left cell for $d=1$, the right
cell for $d=-1$. At its target end the candidate edge is always attached
to state vertex $q$. At its source end attach it to state vertex $p$
if that cell's scanned symbol makes $\widehat\delta(p,b)=(q,d)$.
Otherwise attach it to a fresh degree-one vertex. Different unused
slots have different fresh vertices.

Every actual crossing successor edge occurs in its own slot. Every
other slot is a leaf edge attached to a state vertex at its target end.
Keep every actual stay edge inside its cell, with two distinct
incidences when the edge is a loop. The resulting graph is exactly $G$
with additional leaves and leaf edges. Thus it preserves connectivity
between configurations and preserves the tree property of the component
containing $t$.

This asymmetric attachment rule is what ensures locality. A target cell
never tests its neighbor's symbol: its target attachment is unconditional.
Only the source cell chooses between a state vertex and a fresh leaf,
using its own transition row. Omitting candidates that never occur
removes possible dummy leaves, without removing any actual successor.

\subsection{Contour wires and the two marker tests}

An incidence records one endpoint of an edge at its vertex. A loop has
two distinct incidences at the same vertex. At every vertex of positive
degree choose a cyclic order of its incidences. An arrival along one
incidence is followed by departure along the next incidence in this
order. A degree-one vertex consequently reflects the traversal back
along its edge.

\begin{lemma}[Contour of a tree]\label{btc:mach:contour}
On a finite tree with at least one edge, these successor rules form one
cycle on the directed edge traversals, for every choice of cyclic
orders. Opening the reflection joins at two distinct leaves produces
two paths, each joining a tip at one leaf to a tip at the other.
\end{lemma}

\begin{proof}
The rule permutes the directed traversals: every departure has one
preceding arrival. A one-edge tree has a single cycle of length two.
For induction, remove a leaf $z$ from a larger tree and call its neighbor
$v$. Restrict the cyclic order at $v$ to the remaining incidences.
Reinserting $vz$ replaces one step of the smaller contour by a detour
from $v$ to $z$, back to $v$, and then onward. It inserts two directed
traversals into the existing cycle and creates no new cycle. The contour
is therefore a circular wire. Cutting it at two distinct reflection
joins produces two intervals, each with one endpoint at each cut.
\end{proof}

Only inward boundaries of marker cells receive candidate slots. On the
outward side of the left-marker cell, add a test edge from $c_0$ to a
new outer leaf $z_L$. On the outward side of the right-marker cell, add
a test edge from $t$ to another new outer leaf $z_R$. These two test
edges preserve the tree property of the $t$-component. The test leaves
are connected if and only if the original machine accepts.

Each graph edge now receives two lanes, one in each travel direction.
An incidence has an incoming tip and an outgoing tip. For an edge with
incidences $i,j$, one lane joins the outgoing tip at $i$ to the incoming
tip at $j$, and the other joins the outgoing tip at $j$ to the incoming
tip at $i$. This also applies to loops, whose incidences remain distinct.
At every vertex except $z_L,z_R$, join the incoming tip of each incidence
to the outgoing tip of the next incidence in the chosen cyclic order.
The two outer test-leaf joins are left open.

For each candidate slot, its two ports are labeled by absolute travel
directions $\rightarrow,\leftarrow$. The same enumeration is used on
both sides of every cell. The correspondence of ports to local tips is
\[
 \begin{array}{c|cc}
 \text{cell side}&\rightarrow\text{ lane}&\leftarrow\text{ lane}\\ \hline
 \text{left}&\text{incoming}&\text{outgoing}\\
 \text{right}&\text{outgoing}&\text{incoming}
 \end{array}
\]
Consequently equal-index gluing joins an outgoing half-lane in one cell
to the appropriate incoming half-lane in the next.

At each outward marker side, use ports $1,2$ for the two open tips of
its test edge. The same indices remain ordinary candidate-lane indices
on inward and ordinary boundaries. Pair the remaining outward ports as
$(3,4),(5,6),\ldots,(k-1,k)$, disconnected from all graph edges. This
uses every outward port because $k$ is even; for $k=2$ there are no
remaining ports.

Choose incidence orders separately for each symbol type, then reuse
those orders at every occurrence of that symbol. For example, order
incidences by boundary side and candidate index, followed by state-indexed
stay-edge roles and the marker test role. A loop's two roles are distinct.
Every incidence and its attachment are determined by the local symbol,
so this choice requires no information about an adjacent cell.

\subsection{Local contraction and correctness}

Inside one cell, every lane tip belongs to one lane and one join. The
only exposed endpoints are its boundary ports, each used exactly once.
The separately paired unused outward ports have the same degree-one
property. The finite local wiring therefore has degree two at internal
points and degree one at its ports. Each component containing a port
is a path to exactly one other port. Discard closed components and pair
the two endpoints of every remaining path. This defines a perfect
matching on the cell's $2k$ ports, which is the required $\eta(b)$.

All operations used to define this matching depend only on $b$ and the
machine. In particular, an ordinary symbol has the same diagram at its
first, last, and intermediate occurrences. Gluing equal-index ports of
neighboring cell diagrams reconstructs the full two-lane wiring.
Contracting paths before gluing or after gluing yields the same matching
of outer ports, by the path definition of Brauer multiplication. Thus
\eqref{btc:mach:whole-diagram} represents the complete wiring with the two
test-leaf joins open.

If a wire joins a left test port to a right test port, it stays inside
one component of the augmented graph. The separately paired outward
ports cannot participate. Hence $z_L$ and $z_R$ are connected, and the
machine accepts. Conversely, on an accepted input both test leaves lie
in the tree component of $t$. Temporarily close their reflection joins.
Lemma~\ref{btc:mach:contour} gives one circular contour through both joins;
opening them again gives two paths between the two test sides. Both
left test ports therefore meet right test ports. This proves the
acceptance criterion and completes Theorem~\ref{btc:mach:representation}.

For $w$ empty, there are still two distinct adjacent marker cells. Their
inward boundaries glue by precisely the same slot scheme, and both test
edges are present. The accepted component has at least those two edges,
so the tree-contour argument applies without an exceptional case.
Rejecting cycles elsewhere never enter the proof of the converse and
cannot create a wire between disconnected test leaves.

Extend $\eta$ to body words by multiplication, with
$\eta(\varepsilon)=\Id$. Equality $\eta(u)=\eta(v)$ then implies
that $xuy$ and $xvy$ have the same acceptance status for every pair of
body words $x,y$. This follows by adjoining the same marker diagrams
and context diagrams to the two equal products. This context property
is the precise form needed for the syntactic-monoid arguments below.

\subsection{Why the linear order cannot be improved universally}

The following observation concerns arbitrary recognition by Brauer
monoids, so it also applies to the particular test-port representation.
Fixed marker diagrams and a fixed port predicate can always be absorbed
into an accepting subset of the monoid of body-word diagrams.

\begin{theorem}[A linear degree is sometimes necessary]
\label{btc:mach:linear-optimality}
For every $s\ge2$ there is a binary language recognized by an $s$-state
one-way DFA under the usual end-of-input acceptance convention such that
every monoid recognizer
\[
 \rho:\bits^*\longrightarrow\Br{m},\qquad
 L=\rho^{-1}(P),\quad P\subseteq\Br{m},
\]
has $m\ge s$. The same language has an $(s+1)$-state machine in the
accepting-at-any-time convention of this section. Consequently a universal
degree bound of order $o(s)$ is impossible even over a binary alphabet.
\end{theorem}

\begin{proof}
Use state set $[s]$, initial state $1$, and final set $\{1\}$. Let
$0$ act as the cycle $(1\,2\,\cdots\,s)$ and $1$ as the
transposition $(1\,2)$. These permutations generate the full symmetric
group $\mathfrak S_s$: conjugates of the transposition give the adjacent
transpositions. Finite orders ensure that inverses are represented by
positive words as well.

The syntactic monoid of the accepted language is precisely
$\mathfrak S_s$. Equal word permutations give equal acceptance in every
context. Conversely, suppose two permutations disagree at state $i$.
A prefix can carry state $1$ to $i$, and a suffix can carry one of the
two distinct resulting states to $1$ and the other away from $1$.
The corresponding context distinguishes the two words. This proves the
claimed syntactic monoid directly.

Let $S=\rho(\bits^*)$. If two words have equal $\rho$-images, every
context also has equal images and hence equal acceptance. Therefore the
assignment from a word image to its syntactic class is a well-defined
unital surjection $S\twoheadrightarrow\mathfrak S_s$.
Choose an idempotent $e$ of minimum rank above the group identity and
put $r=\rk(e)$. Lemmas~\ref{btc:alg:corner} and~\ref{btc:alg:unit-lifts}
give a reduced corner
\[
 R_e(eSe)\subseteq\Br{r}
 \quad\text{surjecting onto }\mathfrak S_s.
\]
Every target element is a unit, so every element of this reduced corner
is a permutation. Thus it is a subgroup of $\mathfrak S_r$ and
\[
 s!\le |R_e(eSe)|\le r!\le m!.
\]
For $s\ge2$ this forces $m\ge s$.

To realize the witness in the convention used above, keep all $s$
permutation states nonaccepting, move right on body symbols, and add
one accepting state entered at the right marker exactly from state $1$.
The left-marker rule enters the body without changing the permutation
state. Other right-marker rules may be undefined. This uses $s+1$
states, including the new accepting state, and also handles the empty
word. It proves the stated asymptotic necessity. The argument does not
claim that the constant $8$ in the binary upper bound is optimal.
\end{proof}

\begin{writenote}[Checked, 7 October 2026]
Read line by line, with no error found. Rederived by hand: the
normalization~\eqref{btc:mach:normalization} and its acceptance equivalence;
Lemma~\ref{btc:mach:tree} (a component with $v$ vertices and $v-1$ successor
edges); the slot count $J\le(a+2)s+1$, at most one crossing candidate per
state and scanned symbol plus the single candidate $(1,f,f)$; the induction
of Lemma~\ref{btc:mach:contour}; the acceptance criterion; and
Theorem~\ref{btc:mach:linear-optimality} ($s!\le|R_e(eSe)|\le r!\le m!$).
The lane wiring and its locality were read, not rebuilt independently; the
shipped \file{code/sparse\_slots\_audit.py} builds the local diagrams for 800
machines and passed again at the write (Section~\ref{btc:verify:section}).
\end{writenote}
